\section{Universal relational calculus and torsion}
\label{sec:states-filtration}

This section isolates the categorical part of the argument.  The basic objects are compatible interval configurations, their additive observations, finite labelled correspondences, and affine fibres in the prime-power filtration.  The organizing principle is to retain witness objects and regular epimorphisms rather than replace them by chosen sections.

Let $\mathcal C$ be the set of finite subsets $S\subset\N_{>0}$ whose connected components in the path graph have cardinality $2$ or $3$.  Put
\begin{equation}\label{eq:observations}
W(S)=\sum_{n\in S}\frac1n,
\qquad
\nu(S)=|\pi_0(S)|,
\qquad
\overline W(S)=[W(S)]\in A:=\Q/\Z.
\end{equation}
Let
\[
M=A\times\N,
\qquad
0_M=(0,0),
\qquad
\Omega=(\overline W,\nu):\mathcal C\longrightarrow M.
\]
For $n\ge1$, let
\[
D_n\subseteq\mathcal C^n
\]
be the joint-compatibility domain: $(S_1,\ldots,S_n)\in D_n$ when the supports are pairwise disjoint and their literal union is again an element of $\mathcal C$.  On $D_n$ the partial product is literal union.  Hence
\begin{equation}\label{eq:additive-observations}
W\!\left(\bigcup_iS_i\right)=\sum_iW(S_i),
\quad
\nu\!\left(\bigcup_iS_i\right)=\sum_i\nu(S_i),
\quad
\overline W\!\left(\bigcup_iS_i\right)=\sum_i\overline W(S_i).
\end{equation}
Equivalently, for binary compatible unions the square
\begin{equation}\label{eq:observation-square}
\begin{CD}
D_2 @>{\cup}>> \mathcal C\\
@V{(\Omega\pi_1,\Omega\pi_2)}VV @VV{\Omega}V\\
M^2 @>{+}>> M
\end{CD}
\end{equation}
commutes.  Thus $\mathcal C$, with unit $\varnothing$ and compatible union as partial addition, is a partial commutative monoid, and $\Omega$ is additive on its domain of definition.

\begin{remark}[Abstract scope]\label{rem:abstract-labelled-relations}
Everything in the next two subsections is formal for a partial commutative monoid $\mathscr C$ equipped with an additive map $\Omega:\mathscr C\to M$ to a commutative monoid.  The interval model enters only through the verification of the partial-sum axioms and, later, through the arithmetic estimates and the incompatibility graph.
\end{remark}

\subsection{Relations, base change, and witness objects}

We use the standard double category of sets, maps and relations.  Its vertical arrows are maps, its horizontal arrows are relations, and a square
\[
\begin{array}{ccc}
X & \overset{R}{\rightsquigarrow} & Y\\[-1mm]
{\scriptstyle f}\downarrow && \downarrow{\scriptstyle g}\\[-1mm]
X' & \overset{S}{\rightsquigarrow} & Y'
\end{array}
\quad\Downarrow\alpha
\]
means
\[
R\subseteq(f\times g)^{-1}(S).
\]
Horizontal composition is ordinary relational composition.  Every map $f:X\to Y$ has its graph companion
\[
f_* = \{(x,f(x)):x\in X\}:X\rightsquigarrow Y
\]
and converse conjoint $f^*:Y\rightsquigarrow X$; in this sense the double category is the standard relation equipment.  We use only this elementary map--relation calculus and its base-change squares; its horizontal bicategory is denoted $\mathbf{Rel}$.

Existence is kept in the same form.  In $\mathbf{Set}$, and hence in finite sets, a surjection is a regular epimorphism; regular epimorphisms are stable under pullback and finite products, compose, and remain epimorphic after taking a coproduct over a nonempty index set.  Thus a family of witnesses is represented by a map $E\to B$, and ``there is a witness over every $b\in B$'' is recorded as a regular epimorphism
\[
E\twoheadrightarrow B,
\]
not by selecting a section $B\to E$.  This convention will be used throughout the proof.

\subsection{Labelled correspondences and compatible pasting}

Let $X,Y\subseteq M$ be finite.  A \emph{labelled finite correspondence}
\[
R:X\rightsquigarrow Y
\]
consists of a finite edge object $E_R$ and maps
\[
s_R:E_R\to X,
\qquad
\ell_R:E_R\to\mathcal C,
\qquad
t_R:E_R\to Y
\]
satisfying the observable conservation law
\begin{equation}\label{eq:edge-conservation}
t_R(e)=s_R(e)+\Omega(\ell_R(e))
\qquad(e\in E_R).
\end{equation}
Equivalently, the square
\begin{equation}\label{eq:edge-conservation-square}
\begin{CD}
E_R @>{(s_R,\ell_R)}>> X\times\mathcal C\\
@V{t_R}VV @VV{(x,L)\mapsto x+\Omega(L)}V\\
Y @>>> M
\end{CD}
\end{equation}
commutes, where the bottom arrow is inclusion.  The correspondence is \emph{covering} when $t_R:E_R\twoheadrightarrow Y$ is a regular epimorphism.

For correspondences $R:X\rightsquigarrow Y$ and $S:Y\rightsquigarrow Z$, first form the path pullback
\begin{equation}\label{eq:path-pullback}
\begin{CD}
P(S,R) @>{\mathrm{pr}_S}>> E_S\\
@V{\mathrm{pr}_R}VV @VV{s_S}V\\
E_R @>{t_R}>> Y.
\end{CD}
\end{equation}
The pair of label maps gives
\[
(\ell_R\mathrm{pr}_R,\ell_S\mathrm{pr}_S):P(S,R)\longrightarrow\mathcal C^2.
\]
The compatible path object is the pullback
\begin{equation}\label{eq:compatibility-pullback}
\begin{CD}
P^{\mathrm{comp}}(S,R) @>>> D_2\\
@VVV @VVV\\
P(S,R) @>{(\ell_R\mathrm{pr}_R,\ell_S\mathrm{pr}_S)}>> \mathcal C^2.
\end{CD}
\end{equation}
Define $S\odot R$ to have edge object $P^{\mathrm{comp}}(S,R)$, source $s_R\mathrm{pr}_R$, target $t_S\mathrm{pr}_S$, and label obtained from the upper map in \eqref{eq:compatibility-pullback} followed by union $D_2\to\mathcal C$.  The conservation square for $S\odot R$ is the pasted composite of \eqref{eq:observation-square}, \eqref{eq:path-pullback}, and the two conservation squares of $R$ and $S$.

The identity correspondence $1_X$ has edge object $X$, source and target both $1_X$, and constant label $\varnothing$.  A $2$-cell between correspondences with fixed endpoints is a map of edge objects commuting with source, target and label maps.

\begin{lemma}\label{lem:labelled-relation-category}
Compatible pasting makes the finite state sets and labelled correspondences into a bicategory $\mathbf{Corr}_{\mathcal C}^{\Omega}$.  More precisely, both parenthesizations of any finite composite are canonically represented by the same compatible path object
\[
P_n^{\mathrm{comp}}
=
P_n\times_{\mathcal C^n}D_n,
\]
where $P_n$ is the iterated fibre product of the edge objects over the intermediate state sets.  We henceforth suppress these canonical pullback associators.
\end{lemma}

\begin{proof}
The iterated path object $P_n$ is characterized by the universal property of the fibre product, independently of parenthesization.  Pulling it back along $D_n\hookrightarrow\mathcal C^n$ imposes joint compatibility in one step.  On this object every partial union exists and the total label is the literal union of all labels, independent of parenthesization.  The identity statement follows from the universal property of the pullback with the empty label.  The canonical pullback isomorphisms satisfy the usual coherence because they are uniquely determined by their projections.
\end{proof}

The underlying unlabelled relation is the regular image
\begin{equation}\label{eq:underlying-relation}
U(R)=\operatorname{im}(s_R,t_R)\subseteq X\times Y.
\end{equation}
For $E_R\ne\varnothing$ put
\[
\omega(R)=\max\{W(\ell_R(e)):e\in E_R\},
\]
and set $\omega(R)=0$ when $E_R=\varnothing$.

Since \eqref{eq:compatibility-pullback} is a subobject of the unrestricted path pullback, every composable pair has a canonical comparison
\begin{equation}\label{eq:oplax-inclusion}
U(S\odot R)\subseteq U(S)\circ U(R),
\end{equation}
or equivalently a $2$-cell in the relation equipment
\begin{equation}\label{eq:oplax-2cell}
U(S\odot R)
\xRightarrow{\;\vartheta_{S,R}\;}
U(S)\circ U(R).
\end{equation}
Thus
\[
U:\mathbf{Corr}_{\mathcal C}^{\Omega}\longrightarrow\mathbf{Rel}
\]
is a normal oplax homomorphism of bicategories.  Here \emph{oplax} refers to the displayed comparison direction.

For a finite composable string $R_1,\ldots,R_n$, let
\[
\pi_{0,n}:P_n\longrightarrow X_0\times X_n
\]
be the endpoint map and let $\ell:P_n\to\mathcal C^n$ be the label map.  Then
\begin{equation}\label{eq:strictness-endpoint-images}
U(R_n\odot\cdots\odot R_1)
=\operatorname{im}\bigl(\pi_{0,n}|_{P_n^{\mathrm{comp}}}\bigr),
\qquad
U(R_n)\circ\cdots\circ U(R_1)
=\operatorname{im}(\pi_{0,n}).
\end{equation}
Hence the comparison is an identity precisely when
\begin{equation}\label{eq:strictness-endpoint-criterion}
\operatorname{im}\bigl(\pi_{0,n}|_{P_n^{\mathrm{comp}}}\bigr)
=
\operatorname{im}(\pi_{0,n}).
\end{equation}
The information loss is displayed by the commuting image square
\begin{equation}\label{eq:strictness-information-square}
\begin{CD}
P_n^{\mathrm{comp}} @>>> P_n\\
@V{\pi_{0,n}}VV @VV{\pi_{0,n}}V\\
U(R_n\odot\cdots\odot R_1) @>>>
U(R_n)\circ\cdots\circ U(R_1),
\end{CD}
\end{equation}
where the lower objects denote the corresponding regular images.  A stronger witnesswise condition is that the path-label map factor through the compatibility domain:
\begin{equation}\label{eq:compatibility-factorization}
\begin{CD}
P_n @>{\widetilde\ell}>> D_n @>{\iota}>> \mathcal C^n,
\end{CD}
\qquad
\ell=\iota\circ\widetilde\ell.
\end{equation}
By the universal property of \eqref{eq:compatibility-pullback}, this factorization is equivalent to the canonical map
\[
P_n^{\mathrm{comp}}\longrightarrow P_n
\]
being an isomorphism.  It therefore implies
\begin{equation}\label{eq:rel-strictification}
U(R_n\odot\cdots\odot R_1)
=
U(R_n)\circ\cdots\circ U(R_1).
\end{equation}
Moreover, for every composable string,
\begin{equation}\label{eq:mass-subadditive}
\omega(R_n\odot\cdots\odot R_1)
\le\sum_{i=1}^n\omega(R_i),
\end{equation}
because the label of a compatible path is the literal union of its labels and reciprocal value is additive there.

\subsection{Finite torsion and affine fibres in \texorpdfstring{$\Q/\Z$}{Q/Z}}

For $n\ge1$ put
\[
H_n=\{x\in A:nx=0\}.
\]
For a subgroup $H\le A$, let
\[
q_H:A\longrightarrow A/H
\]
be the quotient map, and for $\alpha\in A/H$ write
\[
\operatorname{Fib}_H(\alpha)=q_H^{-1}(\alpha)
\]
for the corresponding affine torsor fibre.  When $H$ is finite, this fibre is a finite coset of $H$.  In particular, every set $\operatorname{Fib}_H(\alpha)\times I$ used below is a finite state object.

\subsection{Affine extension as a pullback of regular epimorphisms}

Let $H\le K\le A$ be finite subgroups and write
\[
\pi_{H,K}:A/H\longrightarrow A/K
\]
for the quotient map.  The subgroup inclusion is encoded by the exact sequence
\begin{equation}\label{eq:affine-extension-exact}
\begin{CD}
0 @>>> K/H @>>> A/H @>{\pi_{H,K}}>> A/K @>>> 0
\end{CD}
\end{equation}
and by the pullback square
\begin{equation}\label{eq:affine-extension-pullback}
\begin{CD}
K @>>> A\\
@V{q_H|_K}VV @VV{q_H}V\\
K/H @>>> A/H.
\end{CD}
\end{equation}
For $\beta\in A/K$ put
\[
\operatorname{Fib}_{\pi_{H,K}}(\beta)
=\pi_{H,K}^{-1}(\beta)\subseteq A/H.
\]
Translation by $\alpha\in A/H$ gives a canonical torsor isomorphism
\begin{equation}\label{eq:fibre-translation-isomorphism}
\operatorname{Fib}_{\pi_{H,K}}(\beta)
\xrightarrow{\;u\mapsto\alpha+u\;}
\operatorname{Fib}_{\pi_{H,K}}\!\bigl(\pi_{H,K}(\alpha)+\beta\bigr).
\end{equation}

\begin{lemma}[Affine torsor extension principle]\label{lem:affine-torsor-extension}
Let $\alpha\in A/H$ and $\beta\in A/K$.  Let $d:Y\to A$ be a map such that $q_Kd$ is constantly $\beta$ and
\begin{equation}\label{eq:affine-increment-regular-epi}
q_Hd:Y\twoheadrightarrow\operatorname{Fib}_{\pi_{H,K}}(\beta)
\end{equation}
is a regular epimorphism.  Put
\[
F^- = \operatorname{Fib}_H(\alpha),
\qquad
F^+ = \operatorname{Fib}_K\!\bigl(\pi_{H,K}(\alpha)+\beta\bigr).
\]
Then the square
\begin{equation}\label{eq:affine-extension-universal-square}
\begin{CD}
F^-\times Y @>{(x,y)\mapsto x+d(y)}>> F^+\\
@V{\mathrm{pr}_Y}VV @VV{q_H}V\\
Y @>{y\mapsto\alpha+q_Hd(y)}>>
\operatorname{Fib}_{\pi_{H,K}}\!\bigl(\pi_{H,K}(\alpha)+\beta\bigr)
\end{CD}
\end{equation}
is a pullback.  Consequently its upper horizontal map is a regular epimorphism.  More generally, if $r:X\twoheadrightarrow F^-$ is any regular epimorphism, then
\[
X\times Y\longrightarrow F^+,
\qquad
(x,y)\longmapsto r(x)+d(y)
\]
is a regular epimorphism.
\end{lemma}

\begin{proof}
The lower arrow in \eqref{eq:affine-extension-universal-square} is the composite of the regular epimorphism \eqref{eq:affine-increment-regular-epi} with the isomorphism \eqref{eq:fibre-translation-isomorphism}.  The square is a pullback: for $y\in Y$ and $z\in F^+$ with
\[
q_H(z)=\alpha+q_Hd(y),
\]
the element $x=z-d(y)$ belongs to $F^-$ and is the unique point mapping to $(y,z)$.  Hence the upper map is the pullback of a regular epimorphism and is therefore a regular epimorphism.  The final assertion follows by composing it with the product regular epimorphism $r\times1_Y$.
\end{proof}

\begin{corollary}[Compatible affine extension]\label{cor:compatible-affine-extension}
Let $\mathscr S$ and $\mathscr B$ be finite parameter sets with label maps
\[
\lambda_{\mathscr S}:\mathscr S\to\mathcal C,
\qquad
\lambda_{\mathscr B}:\mathscr B\to\mathcal C,
\]
and suppose every pair of labels $\lambda_{\mathscr S}(s)$ and $\lambda_{\mathscr B}(b)$ is compatible.  Suppose
\[
\overline W\lambda_{\mathscr S}:\mathscr S\twoheadrightarrow\operatorname{Fib}_H(\alpha)
\]
is a regular epimorphism, $q_K\overline W\lambda_{\mathscr B}$ is constant with value $\beta\in A/K$, and
\begin{equation}\label{eq:compatible-affine-quotient-epi}
q_H\overline W\lambda_{\mathscr B}:\mathscr B\twoheadrightarrow
\operatorname{Fib}_{\pi_{H,K}}(\beta)
\end{equation}
is a regular epimorphism.  Then the compatible-union map
\[
\mathscr S\times\mathscr B\longrightarrow\mathcal C,
\qquad
(s,b)\longmapsto\lambda_{\mathscr S}(s)\cup\lambda_{\mathscr B}(b)
\]
has residue image exactly
\[
\operatorname{Fib}_K\!\bigl(\pi_{H,K}(\alpha)+\beta\bigr).
\]
If $\nu\lambda_{\mathscr B}$ is constant with value $c$, the component-count set is translated by $c$, so its spread is unchanged.
\end{corollary}

\begin{proof}
Apply Lemma~\ref{lem:affine-torsor-extension} to the regular epimorphism $\overline W\lambda_{\mathscr S}$ and the increment map $d=\overline W\lambda_{\mathscr B}$.  Compatibility makes the union map defined on the whole product, and \eqref{eq:observation-square} identifies its residue map with addition.
\end{proof}

\begin{corollary}[Graded affine translation]\label{cor:graded-affine-translation}
Let $I,I'\subseteq\N$ be finite intervals, let $\alpha\in A/H$ and $\beta\in A/K$, and let $\mathscr L$ be a finite label-parameter set with a map
\[
\lambda:\mathscr L\longrightarrow\mathcal C
\]
such that $q_K\overline W\lambda$ is constantly $\beta$.  Put
\[
Y_{I'}=
\{(L,i)\in\mathscr L\times I:i+\nu(\lambda(L))\in I'\}.
\]
Suppose the map
\begin{equation}\label{eq:graded-quotient-regular-epi}
Y_{I'}\longrightarrow
\operatorname{Fib}_{\pi_{H,K}}(\beta)\times I',
\qquad
(L,i)\longmapsto\bigl(q_H\overline W(\lambda(L)),i+\nu(\lambda(L))\bigr)
\end{equation}
is a regular epimorphism.  Put
\[
T^-=\operatorname{Fib}_H(\alpha)\times I,
\qquad
T^+=\operatorname{Fib}_K\!\bigl(\pi_{H,K}(\alpha)+\beta\bigr)\times I'.
\]
Then the edge object
\begin{equation}\label{eq:graded-local-translation}
E_{\mathscr L}(T^-,T^+)
=
\left\{
((x,i),L):
\substack{(x,i)\in T^-,\ L\in\mathscr L,\\
(x+\overline W(\lambda(L)),i+\nu(\lambda(L)))\in T^+}
\right\}
\end{equation}
with label map $((x,i),L)\mapsto\lambda(L)$ and the evident source and target maps is a covering labelled correspondence
\[
\operatorname{Tr}_{\mathscr L}(T^-,T^+):T^-\rightsquigarrow T^+.
\]
\end{corollary}

\begin{proof}
The map in \eqref{eq:graded-quotient-regular-epi} is a regular epimorphism simultaneously over all target counts.  Pulling it back over each point of $I'$ and applying Lemma~\ref{lem:affine-torsor-extension} gives a regular epimorphism onto the corresponding affine $K$-fibre.  Their finite coproduct is the target map of \eqref{eq:graded-local-translation}, so that target map is a regular epimorphism.
\end{proof}

\subsection{The prime-power filtration}

\begin{lemma}\label{lem:torsion-filtration}
For $m,n\ge1$,
\[
H_n=\langle[1/n]\rangle,
\qquad |H_n|=n,
\]
\[
H_m\le H_n\iff m\mid n,
\qquad
H_m\vee H_n=H_{\operatorname{lcm}(m,n)}.
\]
Consequently the nonzero join-irreducible stages are exactly $H_{p^e}$, with $p$ prime and $e\ge1$.
\end{lemma}

\begin{proof}
If $[q]\in H_n$, then $nq\in\Z$, so $[q]=[j/n]$ for some $j\in\Z$; hence $H_n=\langle[1/n]\rangle$ and has order $n$.  The inclusion criterion follows by testing $[1/m]$.  Put $L=\operatorname{lcm}(m,n)$.  Since $\gcd(L/m,L/n)=1$, there exist $a,b\in\Z$ with
\[
aL/m+bL/n=1.
\]
Dividing by $L$ gives $[1/L]\in H_m+H_n$, whence $H_m\vee H_n=H_L$.  The join-irreducibles of the divisibility lattice are precisely the prime powers.
\end{proof}

Accordingly we index the nontrivial steps by prime powers $Q=p^e$.  The \emph{rank} of the stage $H_{p^e}$ means the positive integer $Q=p^e$ itself.  Let $F_Q$ be the join of all prime-power stages of rank at most $Q$, and $F_{<Q}$ the join of those of smaller rank.  The simple quotient
\begin{equation}\label{eq:simple-quotient}
S_Q=F_Q/F_{<Q}
\end{equation}
has order $p$: the $p$-adic valuation of the lower least common multiple is $e-1$, and adjoining $p^e$ raises it by one.  Thus $S_Q$ is canonically a one-dimensional $\F_p$-object.

For $X\ge1$ define the endpoint stage
\begin{equation}\label{eq:endpoint-stage}
E_X=H_{\operatorname{lcm}(1,\ldots,\lfloor X\rfloor)}.
\end{equation}
Every class $[a/b]\in A$ belongs to $H_b$, so
\begin{equation}\label{eq:cofinal-union}
A=\bigcup_{X\ge1}E_X.
\end{equation}
Equivalently, for every $B$ the quotient system satisfies
\begin{equation}\label{eq:centre-kernel-union}
A/E_B
=
\bigcup_{X\ge B}\ker\!\left(\pi_{B,X}:A/E_B\to A/E_X\right).
\end{equation}
Thus every affine centre dies along a final segment of the endpoint filtration without choosing a representative in $A$.

Finally, if $Q$ is a prime-power step, then
\begin{equation}\label{eq:endpoint-current-step}
E_{Q-1}=F_{<Q},
\qquad
E_Q=F_Q.
\end{equation}
Indeed a prime-power rank is $<Q$ exactly when it is at most $Q-1$.  Thus the local simple quotient \eqref{eq:simple-quotient} is literally the associated quotient of the global endpoint filtration at $Q$.
